% ============================================================================
% CAPÍTULO 4 - RAG (RETRIEVAL-AUGMENTED GENERATION)
% ============================================================================

\chapter{RAG: Retrieval-Augmented Generation}\label{cap:rag}

Este capítulo descreve a implementação do sistema \acrfull{rag}, que permite ao assistente fundamentar suas respostas em protocolos clínicos reais, reduzindo o risco de alucinações e aumentando a confiabilidade das informações fornecidas.

\section{Motivação e Arquitetura}

Modelos de linguagem, mesmo após fine-tuning, podem gerar informações imprecisas ou desatualizadas --- um problema crítico no domínio médico. O RAG resolve essa limitação ao recuperar documentos relevantes de uma base de conhecimento antes de gerar a resposta, permitindo que o modelo fundamente suas respostas em fontes verificáveis.

\begin{figure}[H]
\centering
\begin{tikzpicture}[
    node distance=0.6cm and 1cm,
    block/.style={rectangle, draw, rounded corners=3pt, minimum width=2.8cm, minimum height=0.9cm, align=center, font=\small\sffamily},
    process/.style={block, fill=fiap-red!15, draw=fiap-red, thick},
    data/.style={block, fill=blue!10, draw=blue!60, thick},
    io/.style={block, fill=green!10, draw=green!60!black, thick},
    arrow/.style={->, thick, >=stealth}
]

% Fluxo RAG
\node[io] (query) {Pergunta\\do Usuário};
\node[process, right=of query] (embed) {Gerar\\Embedding};
\node[data, right=of embed] (chroma) {ChromaDB\\(544 docs)};
\node[process, right=of chroma] (retrieve) {Recuperar\\Top-K};
\node[process, below=1cm of retrieve] (prompt) {Construir\\Prompt};
\node[process, left=of prompt] (llm) {LLM\\Fine-tuned};
\node[io, left=of llm] (response) {Resposta\\Fundamentada};

% Setas
\draw[arrow] (query) -- (embed);
\draw[arrow] (embed) -- (chroma);
\draw[arrow] (chroma) -- (retrieve);
\draw[arrow] (retrieve) -- (prompt);
\draw[arrow] (query.south) -- ++(0,-0.5) -| (prompt.north);
\draw[arrow] (prompt) -- (llm);
\draw[arrow] (llm) -- (response);

% Label do contexto
\node[below=0.15cm of prompt, font=\tiny\sffamily\itshape] {Query + Contexto};

\end{tikzpicture}
\caption{Fluxo do sistema RAG: a pergunta é convertida em embedding, documentos relevantes são recuperados do ChromaDB, e o contexto é injetado no prompt do LLM.}
\label{fig:fluxo_rag}
\end{figure}

\subsection{Componentes do Sistema}

O sistema RAG implementado possui três componentes principais:

\begin{enumerate}
    \item \textbf{Chunker Semântico:} Divide os protocolos Markdown em chunks que respeitam a estrutura de seções do documento;
    \item \textbf{Indexador:} Gera embeddings e persiste os chunks no ChromaDB;
    \item \textbf{Retriever:} Busca os chunks mais relevantes para uma query e formata o contexto para o LLM.
\end{enumerate}

\section{Modelo de Embeddings}

A escolha do modelo de embeddings é fundamental para a qualidade da recuperação. Como os protocolos estão em português brasileiro, é necessário um modelo multilíngue com bom desempenho em português.

\subsection{Análise de Alternativas}

A Tabela~\ref{tab:modelos_embedding} apresenta os modelos considerados:

\begin{table}[H]
\centering
\caption{Modelos de embeddings avaliados para o sistema RAG.}
\label{tab:modelos_embedding}
\rowcolors{2}{fiap-light}{white}
\begin{tabularx}{\textwidth}{l c c X}
\toprule
\textbf{Modelo} & \textbf{Dim.} & \textbf{PT-BR} & \textbf{Observações} \\
\midrule
\texttt{all-MiniLM-L6-v2} & 384 & Ruim & Focado em inglês, rápido \\
\texttt{all-mpnet-base-v2} & 768 & Ruim & Melhor qualidade, apenas inglês \\
\texttt{paraphrase-multilingual-MiniLM-L12-v2} & 384 & Boa & Multilíngue, compacto \\
\texttt{paraphrase-multilingual-mpnet-base-v2} & 768 & Boa & Multilíngue, mais robusto \\
\rowcolor{green!10}
\texttt{intfloat/multilingual-e5-base} & 768 & Excelente & Estado da arte para retrieval \\
\bottomrule
\end{tabularx}
\end{table}

\subsection{Modelo Selecionado: E5 Multilingual}

O modelo \texttt{intfloat/multilingual-e5-base} foi selecionado por apresentar o melhor desempenho em tarefas de retrieval multilíngue, incluindo português. Este modelo pertence à família E5 (\textit{EmbEddings from bidirEctional Encoder rEpresentations}), desenvolvida pela Microsoft Research.

\begin{table}[H]
\centering
\caption{Características do modelo E5 Multilingual Base.}
\label{tab:e5_specs}
\rowcolors{2}{fiap-light}{white}
\begin{tabular}{l l}
\toprule
\textbf{Característica} & \textbf{Valor} \\
\midrule
Arquitetura base & XLM-RoBERTa \\
Dimensão do embedding & 768 \\
Idiomas suportados & 100+ (incluindo português) \\
Tamanho do modelo & ~1.1 GB \\
Max sequence length & 512 tokens \\
Treinamento & Contrastive learning em pares de retrieval \\
\bottomrule
\end{tabular}
\end{table}

\begin{infobox}[Prefixos do Modelo E5]
O modelo E5 utiliza prefixos específicos para distinguir entre documentos e queries:
\begin{itemize}
    \item \textbf{Documentos:} \texttt{"passage: <texto do documento>"}
    \item \textbf{Queries:} \texttt{"query: <pergunta do usuário>"}
\end{itemize}
Esta distinção melhora significativamente a qualidade da recuperação, pois o modelo foi treinado com essa convenção.
\end{infobox}

\section{Estratégia de Chunking}

O chunking é o processo de dividir documentos longos em fragmentos menores que possam ser indexados e recuperados individualmente. A estratégia de chunking impacta diretamente a qualidade do RAG.

\subsection{Abordagens Consideradas}

\begin{enumerate}
    \item \textbf{Chunking por tamanho fixo:} Divide o texto em chunks de tamanho fixo (ex: 1000 caracteres). Simples, mas pode quebrar o texto no meio de uma frase ou conceito.

    \item \textbf{Chunking por sentenças:} Agrupa sentenças até atingir um limite. Melhor que tamanho fixo, mas não respeita a estrutura do documento.

    \item \textbf{Chunking semântico:} Utiliza a estrutura do documento (seções, parágrafos) para criar chunks com unidades de significado completas.
\end{enumerate}

\subsection{Implementação: Chunking Semântico}

Optamos pelo \textbf{chunking semântico} baseado na estrutura Markdown dos protocolos, que possuem hierarquia clara de seções (\texttt{\#\#}) e subseções (\texttt{\#\#\#}).

\begin{table}[H]
\centering
\caption{Parâmetros do chunking semântico implementado.}
\label{tab:chunking_params}
\rowcolors{2}{fiap-light}{white}
\begin{tabular}{l c l}
\toprule
\textbf{Parâmetro} & \textbf{Valor} & \textbf{Justificativa} \\
\midrule
Divisão primária & Headers MD & Respeita estrutura do protocolo \\
Tamanho máximo & 1500 chars & ~300-400 tokens, contexto suficiente \\
Overlap & 200 chars & Continuidade entre chunks adjacentes \\
Fallback & Por parágrafos & Seções muito grandes são subdivididas \\
\bottomrule
\end{tabular}
\end{table}

O algoritmo de chunking funciona da seguinte forma:

\begin{lstlisting}[caption={Algoritmo de chunking semântico simplificado.}]
def chunk_markdown(texto, max_size=1500, overlap=200):
    # 1. Dividir por headers (## e ###)
    secoes = dividir_por_headers(texto)

    chunks = []
    for secao in secoes:
        if len(secao.conteudo) <= max_size:
            # Secao cabe em um chunk
            chunks.append(secao)
        else:
            # Subdividir secao grande
            sub_chunks = dividir_com_overlap(
                secao.conteudo, max_size, overlap
            )
            chunks.extend(sub_chunks)

    return chunks
\end{lstlisting}

\subsection{Exemplo de Chunking}

A Figura~\ref{fig:exemplo_chunking} ilustra como um protocolo é dividido em chunks:

\begin{figure}[H]
\centering
\begin{tikzpicture}[
    node distance=0.3cm,
    doc/.style={rectangle, draw=fiap-dark, fill=fiap-light, minimum width=10cm, minimum height=0.6cm, font=\small\ttfamily},
    chunk/.style={rectangle, draw=blue!60, fill=blue!10, minimum width=4.5cm, minimum height=0.6cm, font=\small\ttfamily},
    arrow/.style={->, thick, >=stealth, color=fiap-red}
]

% Documento original
\node[doc] (d1) {\#\# Considerações Iniciais};
\node[doc, below=of d1] (d2) {Texto da seção (800 chars)};
\node[doc, below=of d2] (d3) {\#\# Diagnóstico};
\node[doc, below=of d3] (d4) {Texto da seção (2200 chars)};
\node[doc, below=of d4] (d5) {\#\# Tratamento};
\node[doc, below=of d5] (d6) {Texto da seção (1100 chars)};

% Seta
\node[right=1.5cm of d3, font=\Large] (seta) {$\Rightarrow$};

% Chunks
\node[chunk, right=3cm of d1] (c1) {Chunk 1 (800 chars)};
\node[chunk, below=of c1] (c2) {Chunk 2 (1500 chars)};
\node[chunk, below=of c2] (c3) {Chunk 3 (900 chars)};
\node[chunk, below=of c3] (c4) {Chunk 4 (1100 chars)};

% Labels
\node[above=0.2cm of d1, font=\small\sffamily\bfseries] {Protocolo Original};
\node[above=0.2cm of c1, font=\small\sffamily\bfseries] {Chunks Gerados};

% Conexoes
\draw[arrow, dashed] (d2.east) -- (c1.west);
\draw[arrow, dashed] (d4.east) -- (c2.west);
\draw[arrow, dashed] (d4.east) -- (c3.west);
\draw[arrow, dashed] (d6.east) -- (c4.west);

\end{tikzpicture}
\caption{Exemplo de chunking semântico: seções pequenas viram um chunk; seções grandes são subdivididas com overlap.}
\label{fig:exemplo_chunking}
\end{figure}

\section{Indexação no ChromaDB}

O ChromaDB é um banco de dados vetorial open-source, otimizado para armazenamento e busca de embeddings. Foi escolhido por sua simplicidade, persistência local e integração nativa com Python.

\subsection{Processo de Indexação}

O script \texttt{04\_indexar\_protocolos.py} executa a indexação em três etapas:

\begin{enumerate}
    \item \textbf{Carregamento:} Lê os arquivos Markdown das Linhas de Cuidado e protocolos de emergência;
    \item \textbf{Chunking:} Aplica o chunking semântico a cada documento;
    \item \textbf{Indexação:} Gera embeddings e insere no ChromaDB com metadados.
\end{enumerate}

\begin{lstlisting}[style=bash, caption={Executando a indexação de protocolos.}]
# Indexar todos os protocolos (limpa collection existente)
python scripts/04_indexar_protocolos.py --clear

# Saida esperada:
# Linhas de Cuidado: 14 arquivos, 462 chunks
# Protocolos Emergencia: 5 arquivos, 82 chunks
# Total: 544 documentos indexados
\end{lstlisting}

\subsection{Estrutura dos Metadados}

Cada chunk é armazenado com metadados que facilitam a filtragem e rastreabilidade:

\begin{lstlisting}[caption={Estrutura de metadados de um chunk indexado.}]
{
    "id": "Pathway-Hipertensao_3",
    "content": "## Tratamento\n\nOs principios que norteiam...",
    "metadata": {
        "source": "Pathway-Hipertensao",
        "section": "Tratamento",
        "doc_type": "linha_cuidado",
        "chunk_index": 3,
        "total_chunks": 35
    },
    "embedding": [0.023, -0.041, 0.087, ...]  // 768 dimensoes
}
\end{lstlisting}

\subsection{Estatísticas da Indexação}

A Tabela~\ref{tab:stats_indexacao} apresenta as estatísticas finais da indexação:

\begin{table}[H]
\centering
\caption{Estatísticas da indexação no ChromaDB.}
\label{tab:stats_indexacao}
\rowcolors{2}{fiap-light}{white}
\begin{tabular}{l r r r}
\toprule
\textbf{Fonte} & \textbf{Arquivos} & \textbf{Chunks} & \textbf{Média/Arquivo} \\
\midrule
Linhas de Cuidado & 14 & 462 & 33,0 \\
Protocolos de Emergência & 5 & 82 & 16,4 \\
\midrule
\textbf{Total} & \textbf{19} & \textbf{544} & \textbf{28,6} \\
\bottomrule
\end{tabular}
\end{table}

\begin{infobox}[Persistência do ChromaDB]
O ChromaDB persiste os dados em disco no diretório \texttt{chroma\_db/}. Após a indexação inicial, o banco pode ser carregado instantaneamente em execuções subsequentes, sem necessidade de reindexar.
\end{infobox}

\section{Recuperação de Contexto}

O componente \texttt{ProtocolRetriever} é responsável por buscar os chunks mais relevantes para uma pergunta do usuário.

\subsection{Processo de Busca}

\begin{enumerate}
    \item \textbf{Embedding da query:} A pergunta do usuário é convertida em embedding usando o prefixo \texttt{"query: "};
    \item \textbf{Busca por similaridade:} O ChromaDB calcula a distância entre o embedding da query e os embeddings dos documentos;
    \item \textbf{Ranking:} Os K documentos mais próximos são retornados, ordenados por relevância;
    \item \textbf{Formatação:} O contexto é formatado com metadados de fonte para injeção no prompt.
\end{enumerate}

\begin{lstlisting}[caption={Exemplo de uso do ProtocolRetriever.}]
from src.assistant.rag import ProtocolRetriever

retriever = ProtocolRetriever(
    persist_directory="chroma_db",
    embedding_model="intfloat/multilingual-e5-base"
)

# Buscar protocolos relevantes
results = retriever.search(
    query="Como tratar hipertensao resistente?",
    n_results=3
)

# Obter contexto formatado para prompt
context = retriever.get_context_for_query(
    query="Como tratar hipertensao resistente?",
    n_results=3,
    max_context_length=4000
)
\end{lstlisting}

\subsection{Exemplo de Recuperação}

A seguir, um exemplo real de recuperação para a query \textit{``Quais os sintomas de infarto agudo do miocárdio?''}:

\begin{table}[H]
\centering
\caption{Resultados de busca para query sobre IAM.}
\label{tab:exemplo_busca}
\rowcolors{2}{fiap-light}{white}
\small
\begin{tabularx}{\textwidth}{c l l c X}
\toprule
\textbf{\#} & \textbf{Fonte} & \textbf{Seção} & \textbf{Dist.} & \textbf{Trecho} \\
\midrule
1 & iam & Considerações Iniciais & 0.212 & ``O IAM é uma emergência cardiovascular que requer reconhecimento e tratamento imediatos...'' \\
2 & iam & Exames Diagnósticos & 0.298 & ``ECG em até 10 minutos da chegada. IAMCSST: supradesnível de ST...'' \\
3 & Pathway-IC & Diagnóstico & 0.341 & ``Sinais e sintomas de insuficiência cardíaca incluem dispneia...'' \\
\bottomrule
\end{tabularx}
\end{table}

\subsection{Formatação do Contexto}

O contexto recuperado é formatado para injeção no prompt do LLM:

\begin{lstlisting}[caption={Exemplo de contexto formatado para o prompt.}]
[Fonte: iam] [Secao: Consideracoes Iniciais]
O Infarto Agudo do Miocardio (IAM) e uma emergencia
cardiovascular que requer reconhecimento e tratamento
imediatos. O tempo e critico: "tempo e musculo".

---

[Fonte: iam] [Secao: Exames Diagnosticos]
1. ECG em ate 10 minutos da chegada
   - IAMCSST: supradesnivel de ST >=1mm em 2 derivacoes
   - IAMSSST: infradesnivel de ST, inversao de T

---

[Fonte: Pathway-IC] [Secao: Diagnostico]
Sinais e sintomas de insuficiencia cardiaca incluem
dispneia aos esforcos, ortopneia, edema de MMII...
\end{lstlisting}

\section{Integração com o LLM}

O contexto recuperado é integrado ao prompt do LLM através de um template estruturado que instrui o modelo a basear suas respostas nos protocolos fornecidos.

\begin{lstlisting}[caption={Template de prompt com contexto RAG.}]
Voce e um assistente medico especializado. Use APENAS as
informacoes dos protocolos abaixo para responder. Se a
informacao nao estiver nos protocolos, diga que nao tem
informacao suficiente.

PROTOCOLOS RELEVANTES:
{contexto_recuperado}

PERGUNTA DO MEDICO:
{pergunta_usuario}

RESPOSTA (cite as fontes):
\end{lstlisting}

\begin{warningbox}[Limitação do Contexto]
O contexto do LLM é limitado. O parâmetro \texttt{max\_context\_length} (padrão: 4000 caracteres) controla quanto contexto é injetado no prompt. Recuperar muitos documentos pode exceder o limite e degradar a qualidade da resposta.
\end{warningbox}

\section{Validação do Sistema RAG}

Para validar o funcionamento do RAG, foram realizados testes com queries representativas de diferentes domínios clínicos.

\subsection{Métricas de Qualidade}

\begin{itemize}
    \item \textbf{Relevância:} Os documentos recuperados são relevantes para a query?
    \item \textbf{Cobertura:} O sistema recupera documentos de diferentes fontes quando apropriado?
    \item \textbf{Precisão:} Os documentos mais relevantes aparecem no topo do ranking?
\end{itemize}

\subsection{Resultados dos Testes}

\begin{table}[H]
\centering
\caption{Testes de validação do sistema RAG.}
\label{tab:validacao_rag}
\rowcolors{2}{fiap-light}{white}
\small
\begin{tabularx}{\textwidth}{X l c}
\toprule
\textbf{Query de Teste} & \textbf{Fonte Esperada} & \textbf{Top-1 Correto?} \\
\midrule
``Como tratar hipertensão arterial?'' & Pathway-Hipertensao & Sim \\
``Quais medicamentos para diabetes tipo 2?'' & Pathway-Diabetes & Sim \\
``Sintomas de infarto agudo?'' & iam (emergência) & Sim \\
``Protocolo de sepse grave'' & sepse (emergência) & Sim \\
``Tratamento de asma na crise'' & Pathway-Asma & Sim \\
``Manejo de AVC isquêmico'' & avc (emergência) & Sim \\
``Critérios diagnósticos de depressão'' & Pathway-Depressao & Sim \\
``Antibióticos para pneumonia'' & pneumonia (emergência) & Sim \\
\bottomrule
\end{tabularx}
\end{table}

Todos os testes retornaram o documento esperado como resultado principal (Top-1), validando a qualidade do sistema de recuperação.

\section{Considerações de Desempenho}

\subsection{Tempo de Indexação}

\begin{itemize}
    \item \textbf{Carregamento do modelo:} ~30 segundos (primeira execução, com cache ~5s)
    \item \textbf{Geração de embeddings:} ~2-3 segundos por arquivo
    \item \textbf{Tempo total (19 arquivos):} ~3 minutos
\end{itemize}

\subsection{Tempo de Busca}

\begin{itemize}
    \item \textbf{Embedding da query:} ~50ms
    \item \textbf{Busca no ChromaDB:} ~10ms
    \item \textbf{Tempo total por query:} ~60ms
\end{itemize}

O sistema é suficientemente rápido para uso interativo, com latência imperceptível para o usuário.

\begin{infobox}[Escalabilidade]
O ChromaDB utiliza índices HNSW (\textit{Hierarchical Navigable Small World}) para busca aproximada de vizinhos mais próximos, garantindo tempo de busca logarítmico mesmo com milhões de documentos. Para a escala atual (544 documentos), a busca é praticamente instantânea.
\end{infobox}
